\subsection*{A full R2 witness for both nonidentity resolvents}

\statement{B.55}{Proved}{Both quintic resolvents differ from the identity on a specified full lift set}
There is a nonempty, explicitly specified R2 lift set
$\mathscr L\subset\mathcal H$ such that every $\Phi\in\mathscr L$
satisfies
\[
       \mathfrak T_+(\Phi)\ne\Phi,
       \qquad \mathfrak T_-(\Phi)\ne\Phi.                  \tag{B.55.1}
\]
In particular both homomorphisms of B.53 are nonidentity on
their domain, already on its multiplier-one subgroup.
The distinction is detected on the class-four exponent-seven
quotient of the original level-two free group and of each
resolvent's level-two free group. In the degree-three outer
automorphism quotient their values are
\[
             0,\qquad 30q_3,\qquad30q_3,
                     \qquad30q_3\ne0\text{ over }\mathbb F_7,
                                                               \tag{B.55.2}
\]
where $q_3$ is the actual arithmetic insertion class of B.54.
Thus any member of $\mathscr L$ is a full profinite witness
for the first assertion of P5. The construction does not
assert that a member satisfies the pentagon.

\begin{proof}
We first specify a finite quotient which controls all three
comparisons. Permutations in the following construction act
on the left, with the rightmost factor applied first.
Let $S_3$ act regularly on itself, put
\[
 S=(0\,1),\quad R=(0\,1\,2)\quad\text{in }S_3,
 \qquad
 \bar s=(0\,1)(2\,3),\quad\bar r=(0\,2\,4)
                                      \quad\text{in }A_5.
\]
The five blocks of the induced permutation set are
$\{0,1,2,3,4\}\times S_3$. Define permutations of its thirty
letters by
\[
 \begin{split}
 s_*(i,g)&=(\bar s i,S^{\mathbf1_{i=4}}g),\\
 r_+(i,g)&=(\bar r i,R^{\mathbf1_{i=1}}g),\\
 r_-(i,g)&=(\bar r i,R^{\mathbf1_{i=3}}g).
 \end{split}                                               \tag{B.55.3}
\]
They have the required orders two and three. Hence
$s\mapsto s_*,r\mapsto r_+$ and
$s\mapsto s_*,r\mapsto r_-$ define two homomorphisms from
$B=C_2*C_3$. These are the induced representations of the
two maps
\[
       J\xrightarrow{\rho_\pm}B_\pm
                       \longrightarrow S_3,
       \qquad J=q_5^{-1}(\operatorname{Stab}(4)).
                                                               \tag{B.55.4}
\]
Here $\rho_\pm$ fills one of the two order-three factors
and $B_\pm=C_2*C_3$. To see this identification directly,
the quotient graph of groups for $J$ has five edges,
three vertices for the $C_2$-orbits and three for the
$C_3$-orbits, and its underlying graph is a tree.
The only nontrivial vertex groups are the $C_2$ at the
fixed letter $4$ and the two $C_3$ groups at fixed letters
$1,3$. Assigning $S$ to the first and $R,1$ to the latter
two, or $1,R$ in the other order, is exactly (B.55.4).
There is no unspecified free graph generator.
The names $+$ and $-$ may be interchanged when matching
the two geometric roots in B.53; both are treated identically.

The exact enumerations in
\path{code/check_resolvent_control.py} and
\path{code/check_resolvent_colors.py} show that the two
homomorphisms in (B.55.3) have the same kernel $U$ and that
\[
 \begin{gathered}
 Q=B/U,\qquad |Q|=29160=2^3\cdot3^6\cdot5,\qquad
 |\bar\tau|=10,\\
 U\leq F=\ker(B\longrightarrow S_3),\qquad [F:U]=4860.
 \end{gathered}                                             \tag{B.55.5}
\]
The original quotient $B\to S_3$ sends $s,r$ to $S,R$.
It is tracked from the same words, rather than identified
with a different quotient of the same order.
For clarity, the common-kernel assertion is checked without
assuming an isomorphism between two abstract finite groups:
breadth-first enumeration of $\langle s_*,r_+\rangle$
tracks the minus image and the original $S_3$ image of every
word. Every repeated element has the same two tracked images.
The two resulting images both have order $29160$, proving
equality of their kernels.

The kernel $U$ is torsion-free, because each free factor
of $B$ injects into $Q$. Its cusps are $Q/I$, where
$I=\langle\bar\tau\rangle$ has order ten. Thus there are
$2916$ cusps, its free rank is $4861$, and its compactified
surface has genus $973$. These numbers follow respectively
from the cusp cosets, $1+[B:U]/6$, and
$2g-2+2916=[B:U]/6$. In particular B.54 applies at seven.

We also record the exact images under the filled maps.
The subgroup $J/U$ has order $5832$. Schreier generators
for the stabilizer of block $4$ give three generators,
of orders two, three and three. The normal closure in
$J/U$ of either factor killed in (B.55.4) has order $972$;
the quotient has order six and is the tracked $S_3$.
This calculation is independently checked in
\path{code/check_resolvent_pushforward.py}. Therefore
\[
                  \rho_+(U)=F_+,
                  \qquad\rho_-(U)=F_-,                    \tag{B.55.6}
\]
where $F_\pm$ is the full discrete level-two subgroup of
$B_\pm$. Indeed $B_\pm/\rho_\pm(U)$ is the quotient of
$J/U$ by that normal closure, and $\rho_\pm(U)$ is already
contained in the level-two kernel. Equation (B.55.6) is an
equality of actual images, not only their closures.

\paragraph{The three cusp colourings.}
Each of the maps in (B.55.5)--(B.55.6) sends the cusps of
$U$ to the three labelled cusp classes of its target
level-two group. This gives three colourings of $Q/I$.
The original colouring is $Q$-equivariant through its
quotient $S_3$; a positive meridian $g\tau^{10}g^{-1}$
maps to the fifth power of the corresponding original
cusp meridian. For a resolvent, write the same cusp
using a representative in $J$: multiplication of $g$ on
the right by a power of $\tau$ achieves this because
$A_5=A_4\langle\bar\tau\rangle$.
The unique cusp of $J$ has meridian $\tau^5$, which
$\rho_\pm$ maps to the target cusp generator. Consequently
$g\tau^{10}g^{-1}$ maps to the first power of a level-two
target cusp meridian. Each of the three colours occurs
$972$ times, in each colouring.

Choose the two cusps
\[
                         I,\qquad hI,
                         \qquad h=(rs)^2\in Q.             \tag{B.55.7}
\]
They are distinct. In the original colouring both have
colour zero; in the two resolvent colourings their
colours are respectively $(0,1)$ and $(0,2)$.
The complete trace counts are
\[
\begin{array}{c|c|c}
 \text{colouring}&\text{ordered equal-colour counts}
                         &\text{ordered distinct-colour counts}\\ \hline
 \text{original}&N_{00}=N_{11}=N_{22}=9720&0\\
 \text{either resolvent}&0&N_{ij}=4860\quad(i\ne j).
\end{array}                                                 \tag{B.55.8}
\]
Here each $q\in Q$ contributes the ordered pair
$(qI,qhI)$. The computation uses the permutations
(B.55.3), checks constancy on each right $I$-coset,
and enumerates every one of the $29160$ elements.
The exact data, word and counts are retained in
\path{data/resolvent_color_certificate.json}.
These checks are finite proofs of the specified counts;
no search over a noncofinal family is used.

\paragraph{An actual automorphism realizing the trace.}
Insert the boundary-fixed arithmetic automorphism of B.54
in a disc enclosing the two cusps (B.55.7). Fix a
Schur--Zassenhaus complement $Q$ in
$E=\widehat B/O^7(\widehat U)$, and let $H$ be the closed
pro-seven automorphism group generated by its conjugates
of this insertion, as in (B.54.7).
Choose an element of $C_H(Q)$ whose image on
$N_4(U^{(7)})$ is the complete $Q$-trace.
B.54 proves that this lift set is nonempty.
Extend it to $E$, use the identity at every other relative
prime, and apply R2. This defines the nonempty full lift
set $\mathscr L$ in the statement. Every choice in this
set has the same finite trace. Thus the choices in the
higher automorphism coordinates cannot alter the test
which will prove (B.55.1).

We spell out why the compared actions are actually
controlled by this one relative level. The maps
\[
 U^{(7)}\longrightarrow F^{(7)},\qquad
 U^{(7)}\longrightarrow F_+^{(7)},\qquad
 U^{(7)}\longrightarrow F_-^{(7)}                          \tag{B.55.9}
\]
are all surjective. The last two are induced by the
surjections (B.55.6). For the first, $U\trianglelefteq F$
has index $4860$, prime to seven; its closed image in
$F^{(7)}$ is normal with finite index both a power of
seven and a divisor of $4860$, and hence has index one.
Each map kills $O^7(\widehat U)$ by its defining universal
property. Full R2 lifts act on the original target
through restriction to $F$, and on the other two through
restriction to $J$ followed by $\rho_\pm$. Their action
on $Q$ is the identity, so no further conjugating coset
is needed to preserve $J$ or its two named factor
classes. These comparisons commute as actual quotient
maps after choosing the indicated representatives, and
as canonical outer maps independently of those choices.
In particular (B.55.9) transports the actual trace to
the actions of $\Phi,\mathfrak T_+(\Phi),\mathfrak T_-(\Phi)$.

All three are the identity through degree three in the
nilpotent targets. Their first derivations are computed
in degree four from the boundary classes. Although $U$
has positive genus, no assertion about handle coordinates
is omitted: B.54 lifts the entire trace on $U^{(7)}$,
and the target homology in (B.55.9) is generated by its
three cusp classes. Evaluating the transported derivation
on those classes therefore determines it.

\paragraph{The nonzero finite difference.}
Use $X,Y,Z=-X-Y$ for the target cusp homology classes,
and the notation $U_4,V_4,W_4$ of B.54. For the original
colouring, every translated pair in (B.55.8) has equal
images in homology. Substitution of $5X,5X$, or the
corresponding other colour, in the insertion derivation
is zero. Since five is invertible modulo seven, the
resulting derivation on the original $F^{(7)}$ is zero.

For either resolvent, aggregate the ordered counts into
unordered distinct-colour counts. Each of the three
pairs has count $9720$. There are $972$ source cusps of
each target colour. The same certificate separately
counts the incident translated pairs at each of the
$2916$ individual cusps, counting either endpoint.
For either resolvent every cusp has exactly ten
incidences with each of the two other colours and
none with its own. For the original colouring it
has twenty incidences with its own colour and none
with another colour. These are verified integer
counts; no constancy over the integers is inferred
from a derivation over $\mathbb F_7$ or from the
averaged counts alone.

Thus the transported derivation on $X,Y$ is
\[
       a(-10U_4-20V_4,\ 20V_4+10W_4),
                       \qquad a\in\mathbb F_7^\times,     \tag{B.55.10}
\]
with the same nonzero arithmetic scalar as in B.54.
Subtract the inner derivation with values
\[
                   a(-10U_4+10V_4,\ -10V_4+10W_4).
\]
The remainder is $30a(-V_4,V_4)$, namely $30q_3$ in the
outer quotient. B.54 proves $q_3\ne0$, and $30\equiv2$
modulo seven. The calculation is checked independently
in \path{code/check_resolvent_trace.py}.
It takes place in an actual finite class-four group,
so a nonzero value proves inequality of the full outer
automorphisms. Equations (B.55.2) and (B.55.1) follow.
\end{proof}

These full witnesses are outside the Galois image: their
multiplier is one, so an arithmetic witness would lie in
$G_{\mathbb Q(\sqrt{-15})}$, where B.53 makes both
resolvents the identity. They are also outside
$\mathcal K$, since the nonzero target action contradicts
B.53's kernel containment. Neither conclusion determines
whether a witness satisfies the pentagon. The remaining
assertion of P5, that every element of $\widehat{GT}$ in
the domain is fixed, is still Open and is not inferred
from the arithmetic identity.
